\documentclass[10pt,a4paper]{amsart}
\usepackage[margin=19mm]{geometry}
\usepackage{amsmath,amssymb,mathtools}
\usepackage[scr=rsfs,cal=euler]{mathalfa}
\usepackage{xfrac}
\usepackage[unicode,hidelinks,bookmarks=false]{hyperref}
\newcommand{\ie}{i.e.\ }
\newcommand{\cf}{cf.\ }
\newcommand{\resp}{respectively\ }
\newcommand{\Subsection}{\subsection}
\providecommand{\nobreakdash}{\nobreak\hbox{-}}
\setlength{\emergencystretch}{2em}
\multlinegap0pt
\usepackage{bm}
\usepackage{bbm}
\usepackage{dsfont}
\usepackage{xspace}
\usepackage{cancel}
\usepackage{mathtools}
\usepackage{amsthm}
\makeatletter
\@ifundefined{Theorem}{\newtheorem{Theorem}{Theorem}}{}
\@ifundefined{Lemma}{\newtheorem{Lemma}[Theorem]{Lemma}}{}
\makeatother
\makeatletter
\expandafter\ifx\csname customproof\endcsname\relax
\newenvironment{customproof}[1]{%
  \par\pushQED{\qed}%
  \normalfont\topsep6pt \trivlist
  \item[\hskip\labelsep\itshape
    Proof of #1\@addpunct{.}]\ignorespaces
}{%
  \popQED\endtrivlist\@endpefalse
}
\fi
\newcommand{\dd}{\, \mathrm{d}}
\makeatother
\title[Weighted Perimeters and pth Moments of Inertia of Convex Cur]{Weighted Perimeters and pth Moments of Inertia of Convex Curves and Surfaces}
\author{Gyula Csat\'o, Davide Giovagnoli, Prosenjit Roy}
\date{Source: arXiv:2608.19851v1 (CC BY 4.0)}
\begin{document}
\maketitle

\noindent\emph{Source and licence.} Weighted Perimeters and pth Moments of Inertia of Convex Curves and Surfaces. Version arXiv:2608.19851v1 (CC BY 4.0), distributed under the Creative Commons Attribution 4.0 International licence, as recorded in the arXiv licence metadata for this exact version and shown on the abstract page https://arxiv.org/abs/2608.19851v1. Full rights notice: licenses/arxiv-2608.19851-CC-BY-4.0.txt.\par\smallskip

\noindent\textit{Editorial scope.} Counted unit: the Lemma whose source label is \texttt{lemma:first variation of Ep} in the article (source0.tex lines 968--977, source label \texttt{lemma:first variation of Ep}), delivered together with the article's own complete original proof of it (source0.tex lines 1063--1191, one \texttt{proof} environment of 129 lines). No mathematics is written, repaired or re-derived: statement and proof are the article's own text line for line. The proof is detached from the statement in the source: the article states the result at lines 968--977 and prints its proof at lines 1063--1191, and the proof environment's own header names the statement, so the association is the article's own. Required context retained uncounted: none. Every definition, display and label of those lines is retained unchanged. Omitted from this package and not redistributed: 1--967, 978--1062, 1192--1968. Nothing inside a retained excerpt is omitted or altered: every retained body is delivered exactly as the article's lines are written, so no omission receipt of any kind accompanies this package. Citations: the retained excerpts carry no citation command at all, so no bibliography entry of the article is transcribed and no reference is redistributed. Reference adaptation: none was needed and none was made; every reference inside the delivered unit resolves inside it, so no unresolved reference or citation is produced. Printed slips kept literal: no printed word, symbol, hypothesis or number of the counted statement or of its proof was altered. Layout and class: the article's own class is \texttt{amsart} with packages including graphicx, indentfirst, csquotes, babel, fancyhdr, subfig, hyperref, tikz, dsfont, amssymb, mathtools, amsmath, latexsym, amsthm; the wrapper is the lane's pinned \texttt{amsart} editorial wrapper at 10pt on A4, which restates the theorem-like environments the retained text needs and adds the article's own macro definitions that the retained text needs (\texttt{\textbackslash dd}), with extra wrapper packages (bm, bbm, dsfont, xspace, cancel, mathtools). The delivered numbering is the unit's own. Assets: no retained span contains a figure environment, a graphics-inclusion command, a PDF-page inclusion, a table or a TikZ picture; no image file is copied into this package, the package declares no asset dependency and no image file is redistributed with it. Licence: version arXiv:2608.19851v1 of this article is distributed under the Creative Commons Attribution 4.0 International licence, as recorded in the arXiv licence metadata for this exact version and shown on the abstract page https://arxiv.org/abs/2608.19851v1. A full rights notice is delivered at licenses/arxiv-2608.19851-CC-BY-4.0.txt. Characters: every retained excerpt is pure ASCII, so the delivered engine, XeLaTeX with its default Latin Modern text font, reports no missing character.
\begin{Lemma}\label{lemma:first variation of Ep}
Let $n\geq 2$ be an integer and $ p\in(1-n,\infty)\setminus\{0\}$. Then 
\begin{equation} \label{eq:thesis_first}
    E_p' =\begin{cases}
          \infty   & \text{ if } p<2-n, \\
      >0 & \text{ if } p = 2-n, \\
    0 & \text{ if } p>2-n. \\
      \end{cases}
\end{equation}
\end{Lemma}

\begin{proof} We will use the formulas obtained for the first variation in the proof of Lemma \ref{lemma:first variation of Ep}.

The second variation of the energy on the lateral surface gives
\begin{align*}
    \frac{d^2}{dr^2}E_p(M_r) &= (n-2)(n-3)r^{n-4}\int_0^{h(r)}(r^2+x_n^2)^{\frac{p}{2}}\dd x_n  \\
    &+(n-2)r^{n-3}(r^2+h(r)^2)^{\frac{p}{2}}  h'(r)  \\
    &+(n-2)r^{n-3} \int_0^{h(r)}\frac{d}{dr} \left((r^2+x_n^2)^{\frac{p}{2}} \right)\dd x_n  \\
    &+ (n-2)r^{n-3}(h(r)^2+r^2)^{\frac{p}{2}}h'(r)  \\
    &+ r^{n-2}\left[p(h(r)^2+r^2)^{\frac{p}{2}-1} \left( r+h(r)h'(r)\right)h'(r) +(h(r)^2+r^2)^{\frac{p}{2}}h''(r) \right] \\
    &+  p (n-2)\,  r^{n-3} \int_{0}^{h(r)}  r \left( r^2 + t^2 \right)^{\frac{p}{2}-1} \dd t \\
    &+  p \,  r^{n-2} r\left( r^2 + h^2(r)\right)^{\frac{p}{2}-1}h'(r) \\
    &+ p \, r^{n-2} \int_0^{h(r)} \frac{d}{dr} \left( (r^2 +t^2)^{\frac{p}{2}-1} r \right) \dd t.
\end{align*}

When plugging $r=1$ all the terms having the integral vanish since $h(1)=0$ and we are left with
\begin{align}\label{EpMr double prime}
   \lim_{r\nearrow 1}  \frac{d^2}{dr^2}E_p(M_r) = 2(n-2)h'(1) + 2p h'(1) +h''(1)   =  2-n-2p,
\end{align}
since $h'(1)=-1$ and $h''(1)=n-2$.

Computing the second variation contribution of the bases $D_r$ yields
\begin{align*}
    \frac{d^2}{dr^2}E_p(D_r)&= (n-2) r^{n-3} \left(r^2 +h(r)^2 \right)^{\frac{p}{2}}   \\
    &+ p r^{n-2} \left(r^2 +h(r)^2 \right)^{\frac{p}{2}-1} \left(r+h(r)h'(r) \right)   \\
    &+ p   \left( h(r)h''(r)+h'(r)^2\right) \int_0^r  \left( t^2 +h(r)^2 \right)^{\frac{p}{2}-1} t^{n-2} \dd t  \\
    &+ph(r)h'(r) \left[ \left(r^2+h(r)^2 \right)^{\frac{p}{2}-1} r^{n-2} + (p-2) h(r)h'(r)\int_{0}^{r} \left( t^2 +h(r)^2 \right)^{\frac{p}{2}-2} t^{n-2} \dd t\right].
\end{align*}
Using that $h(1)=0,$ $h'(1)=-1,$ and $h''(1)=n-2$ we can write the second derivative of $E_p(D_r)$ as
\begin{align}\label{EpDr double prime with limit}
    \lim_{r\nearrow 1}\frac{d^2}{dr^2}E_p(D_r) &= \lim_{r\nearrow 1} \left\{ n-2 + p+p [(n-2)h(r)J_1(r)+J_1(r) + (p-2)J_2(r)]\right\},  
     \end{align}
where
\begin{align*}
    J_1(r)&:= \int_0^r  \left( t^2 +h(r)^2 \right)^{\frac{p}{2}-1} t^{n-2} \dd t, \\
    J_2(r)&:= h(r)^2 \int_{0}^{r} \left( t^2 +h(r)^2 \right)^{\frac{p}{2}-2} t^{n-2} \dd t.
\end{align*}


To analyze the above expression as $r \nearrow 1$, we consider different ranges of $p$.

\textit{\underline{Case $p>5-n$.}} 
In this case, all terms in \eqref{EpDr double prime with limit} are finite and thus taking the limit gives

     \begin{align*}
  \lim_{r\nearrow 1} \frac{d^2}{dr^2}E_p(D_r) &=   n-2 + p+p  \int_0^{1} t^{p+n-4} \dd t=  n-2+p + \frac{p}{p+n-3}.
 \end{align*}

Thus summing with the contribution \eqref{EpMr double prime} of $M_r$ leads to 

\begin{align}\label{eq:E_p_doublep>5-n}
    E_p'' &=   2-n-2p +   n-2+p + \frac{p}{p+n-3} =-p \; \frac{p+n-4}{p+n-3}.
\end{align}
%%%%%%%%%%%%%%%%%%%%%%%%%%%
\textit{\underline{Case $3-n<p\leq 5-n$.}}  
In this case we have $\lim_{r\nearrow 1}J_1=\frac{1}{p+n-3}$. In order to deal with $J_2$ we use the change of variables $t=h(r)s$ and obtain
\begin{align}\label{eq:J_2_contr}
    J_2(r)=h(r)^{n+p-3} \int_0^{\frac{r}{h(r)}} s^{n-2} (1+s^2)^{\frac{p}{2}-2} \dd s.
\end{align}
Whenever $p <5-n$ the integral is finite and thus $\lim_{r \nearrow 1} J_2=0$. In the borderline case $p=5-n$ for $s>1$ we have
\[
\int_{1}^{\frac{r}{h(r)}} s^{n-2} (1+s^2)^{\frac{p}{2}-2} \dd s \leq \int_{1}^{\frac{r}{h(r)}} s^{-1} \dd s  = \log \left( \frac{r}{h(r)}\right).
\]
 Computing $\lim_{r \nearrow 1} h(r)^2  \log \left( \frac{r}{h(r)}\right)$ reveals that $\lim_{r \nearrow 1} J_2$ vanishes.
 
 This implies that also in this range we obtain \eqref{eq:E_p_doublep>5-n}.



\smallskip  
\textit{\underline{Case $p=3-n$.}}
In this case, we analyze the exact explosive behavior of $J_1(r)$. Indeed, from \eqref{eq:J_2_contr} we deduce that the term $J_2(r)$ gives a finite contribution. For $J_1(r)$, using once again $t=h(r)s$, we obtain
\begin{align} \label{eq:J_1}
    J_1(r) = h(r)^{p+n-3} \int_0^{\frac{r}{h(r)}} (s^2+1)^{\frac{p}{2}-1} s^{n-2} \dd s = \int_0^{\frac{r}{h(r)}} (s^2+1)^{\frac{1-n}{2}} s^{n-2} \dd s.
\end{align}

As $r\nearrow 1$, the integrand behaves asymptotically as $s^{-1}$. Thus, $J_1(r)$ diverges  as $\log(r/h(r))$.
On the other hand $h(r)J_1(r)$ vanishes as $r \nearrow 1$.
Therefore, since the contribution of the lateral surface is finite, \eqref{EpDr double prime with limit} shows that the behavior of $E_p''$ is uniquely determined by the unbounded term $p J_1(r)$.
For $n>3$, since $p=3-n<0$, we get $E_p''=-\infty$.
While if $n = 2$, then $p = 1$ and we obtain that $E_p'' = \infty$.


\smallskip
\textit{\underline{Case $2-n<p<3-n$.}} 
In the remaining case where $2-n < p < 3-n $  we have that both $J_1$ and $J_2$ are unbounded as $r \nearrow 1$. For this, combining \eqref{EpMr double prime} and \eqref{EpDr double prime with limit}, we express $E_p''$ as
\begin{equation} \label{asymptotic_d2E}
 E_p'' = -p+p[(n-2)h(r)J_1(r)+J_1(r)+(p-2)J_2(r)] +  o(1)\quad \text{as }r \nearrow 1.
\end{equation}
To deal with the case $p \in (2-n,3-n)$ we express conveniently $J_1(r) + (p-2)J_2(r)$. For this we note that (using simply $n+p-3=(n-1)+(p-2)$ in the second identity)
\begin{align*}
	\frac{\dd}{\dd t} \left[ t^{n-1} (t^2 + h^2(r))^{\frac{p}{2}-1}\right] &= (n-1)t^{n-2} (t^2 + h^2(r))^{\frac{p}{2}-1} + (p-2)t^{n}(t^2 + h^2(r))^{\frac{p}{2}-2} \\
	&=(n+p-3) t^{n-2}(t^2 + h^2(r))^{\frac{p}{2}-1} - (p-2) h^2(r) t^{n-2}(t^2 + h^2(r))^{\frac{p}{2}-2}.
\end{align*}
Note that both terms on the right hand side are integrable in $t$ on the interval $(0,r)$ for every $r<1$ for $n\geq 2.$ Thus, 
integrating both sides of the above expression with respect to $t$ from $0$ to $r$ yields the relation
\[
r^{n-1}(r^2 + h^2(r))^{\frac{p}{2}-1} =(n+p-3) J_1(r) -(p-2)J_2(r)
=(n+p-2)J_1(r)-J_1(r)-(p-2)J_2(r),
\]
which implies
\[
J_1(r) +(p-2)J_2(r) = (n+p-2)J_1(r) -r^{n-1}(r^2 + h^2(r))^{\frac{p}{2}-1} 
\]

Thus, as $r \nearrow 1$, we can further simplify \eqref{asymptotic_d2E} as
\begin{align} \label{eq:final1}
 E_p'' = -p+p\left[(n-2)h(r)J_1(r)+(n+p-2)J_1(r) -r^{n-1}(r^2 + h^2(r))^{\frac{p}{2}-1} \right] + o(1)
\end{align}
as $r \nearrow 1.$ The last term in the brackets $[\ldots]$ tends to $1$ as $r \nearrow 1$.  For $n\neq 2$, using that $p<3-n$ we have that
$$
  \lim_{r\nearrow 1}J_1(r)=\infty.
$$
Similarly, for the first term in \eqref{eq:final1}, since $p>2-n$, we get as in \eqref{eq:J_1}
\[
\lim_{r\nearrow 1} h(r) J_1(r) = \lim_{r\nearrow 1}  h^{n+p-2}(r) \int_0^{\frac{r}{h(r)}} s^{n-2} (1+s^2)^{\frac{p}{2}-1} \dd s =0.
\]
Hence, since $p < 0$ and $(n+p-2) > 0$, we conclude
\[
E_p''= -\infty\qquad \text{ if }p\in (2-n,3-n)\quad\text{ and }p\neq 2.
\]

For $n=2$ and $0<p<1$ the first term in \eqref{eq:final1} vanishes.  In this case $ E_p''$  has the same behavior of $
\lim_{r \nearrow 1}p^2J_1(r)$.
Since $0<p<1$ we obtain $E_p''=\infty$ due to
\[
\lim_{r \nearrow 1}J_1(r) = \lim_{r\nearrow 1}  h^{p-1}(r) \int_0^{\frac{r}{h(r)}} s^{n-2} (1+s^2)^{\frac{p}{2}-1} \dd s =\infty,
\]
which concludes the proof in this case. 
\end{proof}

\end{document}
